% ============================================================================
% Chapter 19: Operating Systems in Depth
% ============================================================================

\chapter{Operating Systems in Depth}
\label{ch:os-depth}

\begin{objectives}
  \item Compare different operating systems (Windows, macOS, Linux, Android, iOS)
  \item Understand multitasking and how the OS manages multiple programs
  \item Learn about user accounts and permissions
  \item Explore the concept of a graphical user interface (GUI)
\end{objectives}

\bytetip[tip]{You already know \emph{what} an OS does. Now let's understand \emph{how} it does it --- and explore the world beyond Windows!}

\section{Comparing Operating Systems}
\label{sec:compare-os}
\index{operating system!comparison}

\begin{table}[H]
\centering
\caption{Popular operating systems compared}
\label{tab:os-compare}
\renewcommand{\arraystretch}{1.3}
\begin{tabularx}{\textwidth}{l l X l}
\toprule
\rowcolor{HeaderRow}
\textcolor{white}{\textbf{OS}} & \textcolor{white}{\textbf{Made By}} & \textcolor{white}{\textbf{Used On}} & \textcolor{white}{\textbf{Cost}} \\
\midrule
\rowcolor{RowLight} Windows & Microsoft & Desktops, laptops & Paid \\
\rowcolor{RowDark} macOS & Apple & MacBooks, iMacs & Free (with Mac) \\
\rowcolor{RowLight} Linux & Community & Servers, desktops & Free \\
\rowcolor{RowDark} Android & Google & Smartphones, tablets & Free \\
\rowcolor{RowLight} iOS & Apple & iPhones, iPads & Free (with device) \\
\rowcolor{RowDark} ChromeOS & Google & Chromebooks & Free (with device) \\
\bottomrule
\end{tabularx}
\end{table}

\section{Multitasking}
\label{sec:multitasking}
\index{multitasking}

\hl{Multitasking} means running multiple programs at the same time. The OS manages this by rapidly switching between programs --- so fast that it feels like everything is running simultaneously.

\begin{center}
\begin{tikzpicture}[node distance=0.5cm]
  \node[process, fill=FunSky!15, draw=FunSky, minimum width=2.2cm] (p1) {Chrome};
  \node[process, fill=FunPurple!15, draw=FunPurple, minimum width=2.2cm, right=of p1] (p2) {Word};
  \node[process, fill=LabGreen!15, draw=LabGreen, minimum width=2.2cm, right=of p2] (p3) {Music};
  \node[process, fill=FunCoral!15, draw=FunCoral, minimum width=2.2cm, right=of p3] (p4) {Paint};
  
  % CPU below
  \node[hwbox, fill=NavyBlue!10, draw=NavyBlue, minimum width=3cm, below=1.5cm of p2.south east] (cpu) 
    {\faMicrochip\ CPU};
  
  \foreach \p in {p1, p2, p3, p4}{
    \draw[NavyBlue, line width=1pt, ->] (\p.south) -- (cpu.north);
  }
  \node[below=0.3cm of cpu, font=\scriptsize\sffamily\itshape, text=NavyBlue] 
    {CPU switches between all programs incredibly fast};
\end{tikzpicture}
\end{center}

\section{GUI --- Graphical User Interface}
\label{sec:gui}
\index{GUI}

A \hl{GUI} (Graphical User Interface)\index{GUI} lets you interact with the computer using visual elements --- icons, buttons, windows, and menus --- instead of typing text commands. All modern operating systems use a GUI.

\bytetip[fun]{Before GUIs existed, people used a \textbf{CLI} (Command Line Interface) --- a black screen where you typed commands. It looked like something from a hacker movie! Some IT professionals still prefer CLI because it's faster for certain tasks.}

\begin{funfact}
The first GUI was developed at \textbf{Xerox PARC} in 1973 --- yes, the same Xerox that makes photocopiers! Steve Jobs visited their lab, was inspired, and created the \textbf{Apple Lisa} and later the \textbf{Macintosh} (1984), which brought GUIs to regular people.
\end{funfact}

\begin{reviewqs}
  \item Name four popular operating systems and who makes each one.
  \item What is multitasking? How does the CPU handle it?
  \item What does GUI stand for and why is it important?
  \item Which operating system is free and open-source?
  \item What is the difference between GUI and CLI?
\end{reviewqs}

\begin{challenge}
If you have access to a tablet or smartphone, compare its interface to a Windows desktop. List \textbf{five differences} and \textbf{five similarities} between the two operating systems' interfaces.
\end{challenge}
